\section{Discussion}
A planner that trusts a source-calibrated region on a new dataset trusts a region that is too small. At $\alpha=0.1$
the confirmatory models lose \gapRawConf, \gapRawConfRev, \wmGapAvW\ and \wmGapNsW\ points of coverage on the four
transfers we tested, and the loss is silent: nothing observable without labels signals it. Four results bear
directly on practice.

\emph{Label-free fixes are not safe defaults.} Reweighting made every forward gap larger, and the normalised score
helped in one direction and hurt in the other.

\emph{Auditing is cheap; repairing is not.} About $10^3$ labelled target agents suffice to detect a gap of this size
with high power. Repairing it with a guarantee that holds under scene clustering needs whole labelled scenes: roughly
\recMMin\ to \recMMax\ of them before the region is both finite and within twice the oracle area. Below that budget the
correct behaviour is to abstain, which the scene-level rule does. Fig.~\ref{fig:protocol} summarises the resulting
procedure.

\emph{A learned predictor of scene loss does not replace labels.} A classifier trained on source labels to predict
whether an agent is missed explained at most \predBestRsq\% of the variance of the true per-scene miss fraction on any
target, and also little on held-out source scenes, so it cannot stand in for target labels, for example as a control
variate that shortens the label budget.

\emph{The unit of labelling matters, and scene size alone does not predict the failure.} On nuScenes (mean
\nsAgentsMean\ agents per scene), agent-level guarantees were nearly useless for AutoBot (\agentCorNsForty\% valid at
$m=40$) but mostly valid for Wayformer on the identical scenes (\agentCorWfForty\%, Sec.~\ref{sec:wf}), tracking how
badly each model's source threshold was miscalibrated there (\rSrcAutobot\% vs.\ \rSrcWayformer\% missed), not the
scene structure, which was the same for both. Agent-level guarantees held on AV2 and on Waymo where scenes are small
and the source threshold less wrong, but not for nuScenes$\to$Waymo (\agentCorNsWmThreeHundred\% valid at $m=300$),
whose source threshold misses \rSrcNsWm\% of agents. Without labels a practitioner cannot tell which case applies, so the scene-level guarantee should be the default.

\begin{figure}[t]
\centering
\fbox{\begin{minipage}{0.95\columnwidth}\small
\begin{enumerate}\setlength{\itemsep}{0pt}
\item \emph{Audit.} Binomial test of coverage $\ge1-\alpha-\epsilon$ on $k\approx10^3$ agents (Sec.~\ref{sec:audit}).
      If not rejected, keep $\hat q_{\mathrm{src}}$ and keep auditing.
\item \emph{Certify.} Test $H:R(\hat q_{\mathrm{src}})>\alpha$ on the $m$ scenes at level $\delta/2$ with the
      betting $p$-value; if rejected, keep $\hat q_{\mathrm{src}}$.
\item \emph{Recalibrate.} Otherwise run scene-level Learn-then-Test on the fixed grid (Sec.~\ref{sec:scene}) at level
      $\delta/2$ and return its threshold.
\item \emph{Abstain.} If no finite threshold is certified, return $+\infty$ and label more scenes (expect this below
      about \recMMin\ scenes).
\end{enumerate}
\end{minipage}}
\caption{Audit-then-repair protocol for a transferred conformal region. Steps 2--4 give $\Pr[R(\hat q)\le\alpha]\ge1-\delta$
(Proposition~\ref{prop:scene}); step 1 only detects.}
\label{fig:protocol}
\end{figure}

\paragraph{What a planner would see} Descriptive, agent-weighted numbers (no guarantee covers them): at the source
threshold, \missAgentsAvNs\% of nuScenes agents and \missAgentsAvWm\% of Waymo agents fall outside the AV2-calibrated
region, against a nominal 10\%, and the shortfall is concentrated in turns: \missRightSg\% of right-turning agents in
Singapore and \missRightBos\% in Boston, against \missStraightBos\% of straight-driving agents in Boston. The tube is
\qSrcAvNs\,m wide where \qOracleAvNs\,m is needed; scene-level repair restores \recQAvNs\,m.

\paragraph{Mechanism} Two observations narrow the cause. For AutoBot the measured input differences explain about a
quarter of the gap, and the loss concentrates in turns under left-hand traffic, which AV2 does not contain; for that
model both point to a change in the conditional behaviour of road users rather than in the scenes themselves,
consistent with Proposition~\ref{prop:ident}. This is an interpretation, not a measurement. The turn pattern
replicates on Wayformer, but the input-level explanation does not: its injected input differences reproduce the whole,
smaller gap (Sec.~\ref{sec:wf}), so how much of the loss is protocol and how much is behaviour depends on the model.
Across AutoBot training budgets the more accurate model lost more coverage, but a second training seed already breaks
this pattern, and so does Wayformer, with accuracy comparable to the confirmatory AutoBot model but under half its
nuScenes gap. Accuracy is therefore not a general predictor of gap size, only a pattern across training budgets of one
architecture.

\paragraph{Limitations} (i) Two architectures, one trained in both directions; Wayformer confirms the direction,
reweighting, turn and repair findings but was only trained forward, so its reverse direction is untested, and its
injection result rests on reconstructed kinematics. (ii) For AutoBot the measured input differences leave most of the
gap unexplained, and unmeasured causes cannot be ruled out. (iii) Coverage is marginal over agents or scenes, not
conditional on a scenario; the manoeuvre breakdown is a diagnostic, not a guarantee. (iv) Proposition~\ref{prop:ident}
covers scale shifts only. (v) The repair budget is measured on three datasets and may differ on others; the registered
evaluation design (disjoint pools, exact evaluation) should be reused rather than the numbers. (vi) Seeds: a second
training seed of the confirmatory AutoBot model gave an AV2$\to$nuScenes loss of \seedGapOne\ points (95\% interval
\seedCiLo--\seedCiHi) against \seedGapZero\ (\seedZeroCiLo--\seedZeroCiHi) for the first, meeting our registered
criterion (same sign, within a factor of two); the intervals overlap, but seed-to-seed variation is as large as the
sampling uncertainty of a single run. A second Wayformer seed (accuracy \wfSeedAde) gave \wfSeedGapOne\ points on
nuScenes against \wfSeedGapZero\ and \wfSeedWmOne\ against \wfSeedWmZero\ on Waymo, so its smaller nuScenes gap is not a
seed artefact. The reverse direction and the injection and recalibration analyses were not re-seeded.

\section{Conclusion}
A conformal guarantee for trajectory prediction does not transfer across driving datasets. On three datasets, two
architectures and six model--target transfers, a forecaster calibrated at 90\% covered between \covTgtMin\ and
\covTgtMax\% of target futures, every time under-covering, never over. Reweighting does not repair the loss and no
label-free method can in general. More accurate AutoBot models lost more coverage across training budgets, but not
between seeds or across architectures, so accuracy is not a general predictor of how badly the guarantee fails.
Labelled target data can repair it: a scene-level Learn-then-Test rule with a correctly capped betting $p$-value
restores the guarantee at every budget we tested, on both architectures, and gives a usable region from about
\recMMin\ labelled scenes. We recommend that any transferred conformal region be audited on target labels and, where it
fails, recalibrated at the scene level (Fig.~\ref{fig:protocol}) before it is used for planning.
